\documentclass[11pt,a4paper]{article}
\usepackage[T1]{fontenc}
\usepackage[utf8]{inputenc}
\usepackage{amsmath,amssymb,amsthm}
\usepackage[margin=1in]{geometry}
\usepackage[colorlinks=true,linkcolor=blue,urlcolor=blue]{hyperref}
\theoremstyle{plain}
\newtheorem{theorem}{Theorem}
\newtheorem{lemma}[theorem]{Lemma}
\newtheorem{proposition}[theorem]{Proposition}
\newtheorem{corollary}[theorem]{Corollary}
\theoremstyle{definition}
\newtheorem{remark}[theorem]{Remark}
% A quoted conjecture can be a single hundred-term formula whose terms are thirty-digit
% numbers. TeX cannot hyphenate those, so without a little stretch every such quote runs off
% the page; the linked-a-file papers are all of that shape.
\setlength{\emergencystretch}{3em}

\title{The empirical recurrence for OEIS A251487, proved by transfer matrix}
\author{Adrian Perez Fontelles\\ \small Independent researcher}
\date{14 September 2026}
\begin{document}
\maketitle

\begin{abstract}
OEIS A251487 counts $(n+1)\times5$ arrays over $\{0,\dots,3\}$ subject to a condition imposed on
every $2\times2$ subblock, and carries an empirical recurrence of order $48$ contributed
by R.~H.~Hardin. It is true, and it is not an empirical matter at all. The condition is local
to each subblock, so the rows of an admissible array are the vertices of a finite digraph
on $S=1024$ vertices whose edges are the admissible consecutive pairs, and the count is a
number of walks: $a(n)=\mathbf{1}^{\!\top}M^{n}\mathbf{1}$. Writing $q$
for the characteristic polynomial of the conjectured recurrence, the recurrence holds for all
$n$ exactly when $\mathbf{1}^{\!\top}M^{j}q(M)\mathbf{1}=0$ for every $j\ge0$, and the
Cayley--Hamilton theorem bounds that to $j<S$. It is therefore a finite computation in exact
integer arithmetic, and it comes out zero.
\end{abstract}

\noindent\small 2020 Mathematics Subject Classification. 05A15, 05B45, 15B36.\normalsize

\section{The sequence and the conjecture}

OEIS A251487 is ``Number of (n+1)X(4+1) 0..3 arrays with no 2X2 subblock having the sum of its diagonal elements less than the minimum of its antidiagonal elements.''. It has offset $1$ and begins
\[
740096, 541815424, 396684056960, 290403118046224, 212595602301764864, 155634922739125247648, 113935706529820385852096, 83408950603225147662340176,\ \dots
\]
The entry carries the following comment:
\begin{quote}\small
Empirical recurrence of order 48 (see link above)

Empirical: a(n) = 1024*a(n-1) - 240642*a(n-2) + 20585760*a(n-3) - 644229599*a(n-4) + 4030343200*a(n-5) + 249555532668*a(n-6) - 4978994676928*a(n-7) + 129232348438541*a(n-8) - 5967182688888128*a(n-9) - 207963148292868298*a(n-10) + 2278691172469602016*a(n-11) + 70525378519869926523*a(n-12) - 277965467682574305888*a(n-13) - 6274268544359385844184*a(n-14) + 41219434287675392752256*a(n-15) + 287959027100733314916452*a(n-16) - 2072160289109676467611456*a(n-17) - 11789823459460401467968408*a(n-18) + 64424631338826073539532800*a(n-19) + 150591431215986466695927388*a(n-20) - 772958707368872247278334464*a(n-21) - 1236803645210469424062457608*a(n-22) + 2386305156498633048224218240*a(n-23) - 82062094967830029115561528288*a(n-24) + 73545580799177848309487091456*a(n-25) - 608069475711925910213753958272*a(n-26) - 2235236611374938480797838458368*a(n-27) - 674129525308578024325070194368*a(n-28) + 5766035924488272967276136322048*a(n-29) + 41367117095438812392075112508544*a(n-30) + 10256266083515680606465417162752*a(n-31) - 33157828046425965627313607228928*a(n-32) - 17854537430914625978387863732224*a(n-33) + 61243952051486860092987351146496*a(n-34) - 65791945790192288483707678064640*a(n-35) - 126793683979723609733270448783360*a(n-36) + 67234132165202331600972594806784*a(n-37) + 33798956081283656059949703659520*a(n-38) - 76873855237472964847051106942976*a(n-39) - 5267302444102393505323376050176*a(n-40) - 53177623269325610925340753920*a(n-41) - 14326435114777969085042003017728*a(n-42) + 496365852062724311273063841792*a(n-43) + 1502965646727719856373374124032*a(n-44) + 14398129550933450359987765248*a(n-45) + 25112777411510178447894773760*a(n-46) + 265038183974773738753228800*a(n-47) + 630508117559336427847680000*a(n-48)
\end{quote}
As of the ``Last modified'' line on the live entry (Jul 23 2025 13:25:32, revision 6) the statement is
still recorded as empirical, and nothing on the entry records it as settled either way.

Written out, the claim is
\begin{equation}\label{eq:conj}
a(n)\;=\;1024\,a(n-1) - 240642\,a(n-2) + 20585760\,a(n-3) - 644229599\,a(n-4) + 4030343200\,a(n-5) + 249555532668\,a(n-6) - 4978994676928\,a(n-7) + 129232348438541\,a(n-8) - 5967182688888128\,a(n-9) - 207963148292868298\,a(n-10) + 2278691172469602016\,a(n-11) + 70525378519869926523\,a(n-12) - 277965467682574305888\,a(n-13) - 6274268544359385844184\,a(n-14) + 41219434287675392752256\,a(n-15) + 287959027100733314916452\,a(n-16) - 2072160289109676467611456\,a(n-17) - 11789823459460401467968408\,a(n-18) + 64424631338826073539532800\,a(n-19) + 150591431215986466695927388\,a(n-20) - 772958707368872247278334464\,a(n-21) - 1236803645210469424062457608\,a(n-22) + 2386305156498633048224218240\,a(n-23) - 82062094967830029115561528288\,a(n-24) + 73545580799177848309487091456\,a(n-25) - 608069475711925910213753958272\,a(n-26) - 2235236611374938480797838458368\,a(n-27) - 674129525308578024325070194368\,a(n-28) + 5766035924488272967276136322048\,a(n-29) + 41367117095438812392075112508544\,a(n-30) + 10256266083515680606465417162752\,a(n-31) - 33157828046425965627313607228928\,a(n-32) - 17854537430914625978387863732224\,a(n-33) + 61243952051486860092987351146496\,a(n-34) - 65791945790192288483707678064640\,a(n-35) - 126793683979723609733270448783360\,a(n-36) + 67234132165202331600972594806784\,a(n-37) + 33798956081283656059949703659520\,a(n-38) - 76873855237472964847051106942976\,a(n-39) - 5267302444102393505323376050176\,a(n-40) - 53177623269325610925340753920\,a(n-41) - 14326435114777969085042003017728\,a(n-42) + 496365852062724311273063841792\,a(n-43) + 1502965646727719856373374124032\,a(n-44) + 14398129550933450359987765248\,a(n-45) + 25112777411510178447894773760\,a(n-46) + 265038183974773738753228800\,a(n-47) + 630508117559336427847680000\,a(n-48)
\end{equation}
for all $n>47$.

\section{The condition is local, so the rows form a digraph}

Write an array of the shape in question as its list of rows
$r^{(1)},r^{(2)},\dots$, each an element of
$V=\{0,1,\dots,3\}^{5}$, so $|V|=S=1024$. A $2\times2$ subblock of the array
occupies two consecutive rows and two consecutive columns; if the two rows are $r$
and $s$ (in that order) and the two columns are numbered $j,j+1$,
the subblock is
\[
\begin{pmatrix}\alpha&\beta\\\gamma&\delta\end{pmatrix}=\begin{pmatrix}r_j&r_{j+1}\\s_j&s_{j+1}\end{pmatrix}.
\]
The entry's requirement on that subblock is
\begin{equation}\label{eq:loc}
\text{not }\bigl((\alpha+\delta)<\min(\beta,\gamma)\bigr),
\end{equation}
a condition on the four entries alone.

For $r,s\in V$ declare $r\to s$ an edge of a digraph $G$ when, for every
$1\le j\le 5-1$,
\begin{itemize}
\item the four entries above satisfy \eqref{eq:loc};

\end{itemize}
Let $M\in\{0,1\}^{S\times S}$ be the adjacency matrix of $G$ and $\mathbf{1}$ the
all-ones vector.

\begin{lemma}\label{lem:walk}
The arrays counted by the entry with $L$ rows are exactly the walks
$r^{(1)}\to r^{(2)}\to\cdots\to r^{(L)}$ of length $L-1$ in $G$. Consequently
\[
a(n)\;=\;\mathbf{1}^{\!\top}M^{n}\mathbf{1} .
\]
\end{lemma}

\begin{proof}
An array with $L$ rows and $5$ columns is precisely a sequence
$r^{(1)},\dots,r^{(L)}$ of elements of $V$; conversely any such sequence is an array.
Every $2\times2$ subblock lies in two consecutive rows $r^{(t)},r^{(t+1)}$ and two
consecutive columns, so the array satisfies the entry's condition on all of its subblocks
if and only if each consecutive pair $\bigl(r^{(t)},r^{(t+1)}\bigr)$ satisfies it for
every column pair --- which is exactly the edge relation of $G$. The number of walks of
length $L-1$, summed over all starting and ending vertices, is
$\mathbf{1}^{\!\top}M^{L-1}\mathbf{1}$, since $(M^{k})_{r,s}$ counts the walks of
length $k$ from $r$ to $s$. The entry's index $n$ corresponds to $L=n+1$ rows.
\end{proof}

\section{The criterion}

Let $q(t)=t^{48}-\sum_i c_i\,t^{48-i}$ be the characteristic polynomial of
\eqref{eq:conj}, with the $c_i$ read off that equation.

\begin{lemma}\label{lem:crit}
For every $n$ with $n+1-1\ge48$,
\[
a(n)-\sum_i c_i\,a(n-i)\;=\;\mathbf{1}^{\!\top}M^{\,n+1-1-48}\,q(M)\,\mathbf{1} .
\]
Hence \eqref{eq:conj} holds from that point on if and only if
$\mathbf{1}^{\!\top}M^{j}q(M)\mathbf{1}=0$ for every $j\ge0$; and it is enough to check
$j=0,1,\dots,S-1$.
\end{lemma}

\begin{proof}
By Lemma~\ref{lem:walk}, $a(n-i)=\mathbf{1}^{\!\top}M^{\,n+1-1-i}\mathbf{1}$, so
\[
a(n)-\sum_i c_i a(n-i)
=\mathbf{1}^{\!\top}\Bigl(M^{m}-\sum_i c_i M^{m-i}\Bigr)\mathbf{1}
=\mathbf{1}^{\!\top}M^{\,m-48}\,q(M)\,\mathbf{1}
\]
with $m=n+1-1$, which is the stated identity; the equivalence follows by letting
$m-48=j$ run over $\ge0$. For the last clause put $w=q(M)\mathbf{1}$. By the
Cayley--Hamilton theorem $M^{S}$ is an integer combination of $I,M,\dots,M^{S-1}$, so
$\mathbf{1}^{\!\top}M^{j}w$ for $j\ge S$ is the corresponding combination of the earlier
values; if those all vanish, so do all the rest.
\end{proof}

\section{The computation}

\begin{theorem}
The recurrence \eqref{eq:conj} holds for every $n>47$, which is the statement of
Section~1.
\end{theorem}

\begin{proof}
The digraph was constructed from \eqref{eq:loc} by a depth-first scan across the $5$
columns: the edge condition is a conjunction of one constraint per consecutive column
pair, so the successors of a row $r$ are enumerated column by column without testing all
$S^2$ pairs. The vector $w=q(M)\mathbf{1}\in\mathbb{Z}^{S}$ was then formed by $48$
matrix--vector products in exact integer arithmetic, and $\mathbf{1}^{\!\top}M^{j}w$ was
evaluated for $j=0,1,\dots$, again exactly. Every one of those integers vanishes from the
index corresponding to $n=47+1$ onwards, and the run continued until $w$ itself was the
zero vector, which settles every larger $j$ at once. By Lemma~\ref{lem:crit} the recurrence
holds for every $n>47$.
\end{proof}

\begin{remark}
The argument gives more than the recurrence: it shows the sequence is $C$-finite with
characteristic polynomial dividing that of $M$, so it has a rational generating function
whose denominator divides $\det(I-xM)$. The conjectured recurrence is one consequence; the
minimal one is a divisor of $q$.
\end{remark}

\section{Verification}

Three checks, each independent of the algebra above.

First, the digraph was built directly from the entry's stated condition and the walk counts
$\mathbf{1}^{\!\top}M^{L-1}\mathbf{1}$ compared against the entry's DATA at
every one of the $9$ published indices; the values agree exactly. That is what ties
the digraph to the sequence: a model that reproduced only some of the published terms would
be the wrong model, and would have been rejected. In particular it is what rules out a
misreading of the entry's wording, since a different reading of the condition gives a
different digraph and different counts.

Second, the recurrence \eqref{eq:conj} was evaluated on those published terms in exact
integer arithmetic, with no matrices involved, and vanishes wherever it applies.

Third, the annihilation test of Lemma~\ref{lem:crit} was run in exact integer arithmetic
until the working vector was identically zero, not to a sampled prefix, and returned zero
throughout.

\begin{thebibliography}{9}
\bibitem{oeis} The OEIS Foundation, \emph{The On-Line Encyclopedia of Integer
Sequences}, \url{https://oeis.org}, 2026. Sequence A251487.
\bibitem{stanley} R.~P.~Stanley, \emph{Enumerative Combinatorics}, Volume 1, 2nd ed.,
Cambridge University Press, 2011. (Transfer-matrix method, Section 4.7.)
\bibitem{fs} P.~Flajolet and R.~Sedgewick, \emph{Analytic Combinatorics}, Cambridge
University Press, 2009.
\bibitem{horn} R.~A.~Horn and C.~R.~Johnson, \emph{Matrix Analysis}, 2nd ed., Cambridge
University Press, 2013.
\end{thebibliography}

\end{document}
